% Осень26!A14; нумерация внутри столбца без пустых строк.
% Источник: лекция 5, §12 (строка 404)
% Источник: лекция 5, §11 (строка 395)
\ticket{12}{Достаточное условие экстремума}
\label{ticket:12}

\begin{definition}
\textbf{Второй дифференциал} функции $f\in C^2$ в точке $a$ --- квадратичная форма
\[
  d^2f(a)[h]=\sum_{i,j=1}^m\partial_i\partial_j f(a)\,h_i h_j.
\]
\end{definition}

\begin{theorem}
Пусть $f\in C^2$ в окрестности $a\in\operatorname{Int}D$, $\operatorname{grad}f(a)=0$, $Q(h)=\tfrac12 d^2f(a)[h]$.
Тогда:
\begin{enumerate}
  \item если $Q$ положительно определена, то $a$ --- строгий локальный минимум;
  \item если $Q$ отрицательно определена, то $a$ --- строгий локальный максимум;
  \item если $Q$ незнакоопределённа, то в $a$ нет локального экстремума.
\end{enumerate}
Если $Q$ имеет постоянный нестрогий знак и обращается в нуль на ненулевом векторе, общего вывода об экстремуме сделать нельзя.
\end{theorem}

\begin{proof}
По формуле Тейлора с остатком Пеано
\[
  f(a+h)-f(a)=Q(h)+o(|h|^2).
\]
Для положительно определённой $Q$ по \hyperref[ticket:11]{лемме об оценке формы} имеем $Q(h)\ge\gamma|h|^2$.
При достаточно малых $h\ne0$ модуль остатка не превосходит $\gamma|h|^2/2$, поэтому
\[
  f(a+h)-f(a)\ge\frac\gamma2|h|^2>0.
\]
Для отрицательно определённой формы применяем рассуждение к $-f$.

Если $Q(u)>0$ и $Q(v)<0$, то
\[
  f(a+tu)-f(a)=t^2Q(u)+o(t^2)>0,
\]
а вдоль $v$ разность отрицательна при достаточно малых $t\ne0$.
Следовательно, экстремума нет.

Последнее утверждение подтверждают $f_1(x,y)=x^2+y^4$ и $f_2(x,y)=x^2-y^4$.
В нуле у них одинаковая квадратичная часть $Q(h)=h_1^2$; у $f_1$ есть минимум, а у $f_2$ экстремума нет.
\end{proof}
